\section{Data}
\label{sec:data}

Every dataset used here is public. Nothing in the analysis requires LIGO
credentials. All availability statements below were verified against the live
services on 2026-08-04 by the environment-check stage of the pipeline
(\cref{app:reproducibility}).

\subsection{Observing-run epochs}

\Cref{tab:epochs} lists the epochs analysed and their GPS boundaries, taken
from the GWOSC data-release pages.

\begin{table}[htbp]
\centering
\caption{Observing-run epochs. Strain from both detectors is public for all
three. O4b and O4c are excluded: their releases were partial at the time of
writing and they add nothing to the O3-versus-O4a comparison.}
\label{tab:epochs}
\small
\begin{tabular}{llrrl}
\toprule
Epoch & Run & GPS start & GPS end & UTC span \\
\midrule
\Otza & O3 & 1238166018 & 1253977218 & 2019-04-01 to 2019-10-01 \\
\Otzb & O3 & 1256655618 & 1269363618 & 2019-11-01 to 2020-03-27 \\
\Ofa  & O4 & 1368975618 & 1389456018 & 2023-05-24 to 2024-01-16 \\
\bottomrule
\end{tabular}
\end{table}

Strain is taken from the \SI{4}{\kilo\hertz} release, which is sufficient for
an analysis confined below \SI{256}{\hertz}. Analysis-ready segments come from
the GWOSC segment API using the \code{<IFO>\_DATA} flag; this is the correct
denominator for a glitch rate, since a rate computed against calendar time
would confound glitchiness with duty cycle.

\subsection{Gravity Spy classifications: O3 only}
\label{sec:data:gravityspy}

The labelled input is the Gravity Spy machine-learning classification release,
Zenodo record 5649212 \citep{GravitySpyData2021,Glanzer2023}. It contains every
Omicron \citep{Robinet2020} trigger with $\mathrm{SNR} > 7.5$ and peak
frequency in \SIrange{10}{2048}{\hertz}, together with the trigger metadata, a
$22$-element vector of class posteriors, and the arg-max label.

Selection cuts, applied identically to the target and control classes so that
both are drawn under the same conditions:

\begin{itemize}
  \item $\mathrm{SNR} \geq 7.5$ (the Gravity Spy trigger threshold);
  \item peak frequency in \SIrange{10}{2048}{\hertz};
  \item machine-learning confidence in the assigned class $\geq 0.90$;
  \item event time inside the epoch boundaries of \cref{tab:epochs}.
\end{itemize}

The confidence cut is the one real choice here. It trades sample size for
label purity, and at $0.90$ it removes the ambiguous tail without collapsing
the sample. \Cref{sec:results:systematics} reports how the results move when
it is varied.

The control sample consists of the classes \emph{Blip}, \emph{Koi\_Fish},
\emph{Low\_Frequency\_Burst}, \emph{Tomte}, \emph{Extremely\_Loud},
\emph{Repeating\_Blips}, \emph{Light\_Modulation}, \emph{Violin\_Mode},
\emph{Whistle}, \emph{Power\_Line} and \emph{No\_Glitch}. Its role is to define
what the classifier of \cref{sec:methods:classifier} must reject.

\paragraph{There is no O4a catalogue.} This is worth stating precisely because
the analysis design turns on it. The Zenodo ML-classification release covers
O1, O2, O3a and O3b. The volunteer-classification release extending to
2024-07-20 is keyed by Zooniverse subject identifier and carries no GPS time,
detector or trigger metadata, so it cannot serve as a trigger list. The O4a
Gravity Spy model of \citet{Wu2025} is published without an accompanying
catalogue. Consequently the O4a sample must be built from strain
(\cref{sec:methods:search}).

\subsection{Auxiliary channels}
\label{sec:data:aux}

Both GWOSC auxiliary releases are served over NDS2 at \code{nds.gwosc.org}.
Only short windows around trigger times, and hour-scale trends, are needed
here, so the several-terabyte bulk downloads are unnecessary.

\subsubsection{O3}

The O3 release \citep{GWOSCO3Aux} comprises $86$ full-rate channels
($40$ \Hone, $46$ \Lone) and $402$ second-trend channels at \SI{1}{\hertz}
($190$ \Hone, $212$ \Lone). \Cref{tab:auxchannels} lists those the analysis
uses.

\begin{table}[htbp]
\centering
\caption{Auxiliary channels used in the O3 analysis. Names ending
\code{.rms} come from the second-trend frames. The suffix
\code{BLRMS\_100M\_300M} denotes a band-limited RMS over
\SIrange{0.1}{0.3}{\hertz}; ``M'' is millihertz.}
\label{tab:auxchannels}
\footnotesize
\begin{tabular}{llp{6.4cm}}
\toprule
Channel & Rate & Role \\
\midrule
\code{L1:SUS-OM3\_M1\_DAMP\_L\_IN1\_DQ}
  & \SI{256}{\hertz}
  & OM3 top-stage \emph{length} damping error signal; the motion witness for
    the per-trigger predictor \\[0.3em]
\code{[HL]1:ISI-GND\_STS\_*\_BLRMS\_100M\_300M.rms}
  & \SI{1}{\hertz}
  & microseismic band at ETMX, ETMY, ITMX, ITMY, HAM5 in $x$, $y$, $z$;
    driver of slow scattering \\[0.3em]
\code{[HL]1:ISI-GND\_STS\_*\_BLRMS\_\{1\_3,3\_10\}.rms}
  & \SI{1}{\hertz}
  & \SIrange{1}{3}{\hertz} and \SIrange{3}{10}{\hertz} anthropogenic bands;
    driver of fast scattering \\[0.3em]
\code{[HL]1:ISI-GND\_STS\_ITMY\_Z\_BLRMS\_30M\_100M}
  & \SI{16}{\hertz}
  & \SIrange{0.03}{0.10}{\hertz} earthquake band \\
\bottomrule
\end{tabular}
\end{table}

Two asymmetries are worth recording. First, \Lone{} has a full-rate suspension
witness in the public release and \Hone{} does not; the per-trigger predictor
of \cref{sec:methods:predictor} is therefore an \Lone-only result, and the
\Hone{} auxiliary analysis is confined to the rate level. Second, the
auxiliary channels are published in raw digital counts with no calibration to
metres, which is dealt with in \cref{sec:methods:predictor}.

\subsubsection{O4a}
\label{sec:data:aux:o4a}

The O4a release \citep{GWOSCO4Aux} is $25$ channels: $14$ from \Hone{} and
$11$ from \Lone. They are reproduced in full below, because the point being
made depends on the list being complete rather than summarised.

\begin{lstlisting}[language=,basicstyle=\ttfamily\scriptsize]
H1: CAL-PCALX_RX_PD_OUT_DQ            L1: ASC-X_TR_A_NSUM_OUT_DQ
    CAL-PCALY_RX_PD_OUT_DQ                CAL-PCALX_RX_PD_OUT_DQ
    IMC-WFS_A_DC_PIT_OUT_DQ               CAL-PCALY_RX_PD_OUT_DQ
    IMC-WFS_A_DC_YAW_OUT_DQ               IMC-WFS_B_I_PIT_OUT_DQ
    LSC-POP_A_LF_OUT_DQ                   OAF-IMC_WFS_B_I_PIT_PREFILT_OUT_DQ
    LSC-REFL_A_RIN_OUT_DQ                 PEM-EX_VMON_ETMX_ESDPOWER24_DQ
    OAF-IMC_WFS_A_DC_PIT_PREFILT_OUT_DQ   PEM-EY_MAINSMON_EBAY_1_DQ
    OAF-IMC_WFS_A_DC_YAW_PREFILT_OUT_DQ   SUS-ETMX_L1_CAL_LINE_OUT_DQ
    OAF-REFL_A_RIN_PREFILT_OUT_DQ         SUS-ETMX_L2_CAL_LINE_OUT_DQ
    PEM-EY_MAINSMON_EBAY_1_DQ             SUS-ETMX_L3_CAL_LINE_OUT_DQ
    SUS-ETMX_L1_CAL_LINE_OUT_DQ           SUS-PI_PROC_COMPUTE_MODE5_RMSMON
    SUS-ETMX_L2_CAL_LINE_OUT_DQ
    SUS-ETMX_L3_CAL_LINE_OUT_DQ
    SUS-PI_PROC_COMPUTE_MODE29_RMSMON
\end{lstlisting}

Every one is a photon-calibrator readout, an input mode-cleaner wavefront
sensor, a length- or angular-sensing readout, a mains voltage monitor, an
electrostatic-drive voltage monitor, a calibration-line drive, or a
parametric-instability RMS monitor.

There is no seismometer, no accelerometer, no microphone and no suspension
witness. Neither the rate-level correlation against ground motion nor the
per-trigger scattering predictor is computable for O4a from public data. This
is a limitation of the release, not of the method, and it is the reason the
O3-versus-O4a comparison in \cref{sec:results:comparison} is a morphology
comparison rather than a mechanism comparison.

The analysis pins this fact in code: a unit test asserts that the O4a channel
list contains none of the tokens \code{ISI-GND}, \code{\_SEIS\_},
\code{\_ACC\_}, \code{SUS-OM} or \code{WIT\_L}. Should GWOSC extend the
release, that test fails, which is the intended signal to revisit the design.

\subsection{Summary of what is and is not possible}

\begin{table}[htbp]
\centering
\caption{Analysis coverage by run, given the public data.}
\label{tab:coverage}
\small
\begin{tabular}{lll}
\toprule
Analysis & O3 & \Ofa \\
\midrule
Q-scan morphology            & yes, on Gravity Spy labels & yes, on classifier output \\
Rate versus ground motion    & yes                        & \textbf{no} --- no seismic channels \\
Per-trigger predictor        & \Lone{} only               & \textbf{no} --- no motion witness \\
Slow/fast decomposition      & yes                        & yes, from arch period alone \\
\bottomrule
\end{tabular}
\end{table}

The slow/fast decomposition survives into O4a precisely because
\cref{eq:fdrive} makes it a strain-only measurement. That is what makes the
comparison possible at all, and it is the reason the auxiliary work in
\cref{sec:results:aux} matters beyond O3: it is the calibration that justifies
reading an O4a arch period as a statement about ground motion.
